\documentclass[10pt]{article}
\usepackage[margin=0.45in,landscape]{geometry}
\usepackage{booktabs}
\usepackage{graphicx}
\usepackage{newtxtext,newtxmath}
\usepackage{microtype}
\pagestyle{plain}
\begin{document}

\noindent\textbf{Table 4: Crisis Prediction with Debt Growth and Real Asset Appreciation by Sector}

\smallskip
\noindent\scriptsize
The estimated country-fixed-effects linear probability model is
\[
 \mathrm{Crisis}_{i,t+1:t+h}
 = \alpha_i^{(h)} + \beta^{(h)}\,\mathrm{High\ Debt\ Growth}_{i,t}
 + \delta^{(h)}\,\mathrm{High\ Price\ Growth}_{i,t}
 + \gamma^{(h)}\,\mathrm{R\mbox{-}Zone}_{i,t} + \varepsilon_{i,t}^{(h)}.
\]
The dependent variable equals one if a crisis occurs within the following
one through four years. High growth indicators equal one in the highest
sector-specific growth quantile; R-Zone is their interaction. Coefficients are
percentage points. Brackets report t-statistics based on Driscoll--Kraay standard
errors with lags 0, 3, 5, and 6. *, **, and *** denote significance at the 10\%,
5\%, and 1\% levels. The sum row reports the model-implied effect when every
included indicator equals one. These are reconstructed-sample estimates and may 
not equal the published paper's values.

\normalsize
\medskip
\input{../_output/table4_replication.tex}

\end{document}
